\DocumentMetadata{
  pdfversion=2.0,
  pdfstandard=ua-2,
  lang=en-US,
  tagging=on,
  tagging-setup={math/setup=mathml-SE,tabsorder=structure}
}
\documentclass[11pt,letterpaper]{article}
\usepackage[margin=0.68in,includefoot]{geometry}
\usepackage{newcomputermodern}
\usepackage{amsmath,amscd,mathtools}
\providecommand{\square}{\mdlgwhtsquare}
\usepackage[hidelinks]{hyperref}
\hypersetup{
  pdftitle={Algebra First-Year Exam, May 2024, Part 2},
  pdfauthor={University of Florida Department of Mathematics},
  pdfsubject={Graduate mathematics examination},
  pdfdisplaydoctitle=true,
  bookmarksopen=true
}
\setcounter{secnumdepth}{0}
\setlength{\parindent}{0pt}
\setlength{\parskip}{0.28em}
\setlength{\emergencystretch}{3em}
\setlength{\leftmargini}{2.15em}
\setlength{\leftmarginii}{2.65em}
\setlength{\leftmarginiii}{3em}
\renewcommand{\labelenumi}{\textbf{\theenumi.}}
\pagestyle{empty}
\raggedbottom
\allowdisplaybreaks
\newenvironment{examproblems}[1]{%
  \begin{enumerate}%
  \setcounter{enumi}{#1}%
  \setlength{\topsep}{0.35em}%
  \setlength{\partopsep}{0pt}%
  \setlength{\itemsep}{0.48em}%
  \setlength{\parsep}{0.18em}%
}{\end{enumerate}}
\newenvironment{examsubparts}{%
  \begin{enumerate}%
  \setlength{\topsep}{0.18em}%
  \setlength{\partopsep}{0pt}%
  \setlength{\itemsep}{0.16em}%
  \setlength{\parsep}{0.1em}%
}{\end{enumerate}}
\newenvironment{examinstructions}{%
  \par\begingroup\itshape
}{%
  \par\endgroup\vspace{0.18em}
}
\makeatletter
\renewcommand\section{\@startsection{section}{1}{0pt}%
  {-1.25ex plus -.25ex minus -.1ex}%
  {0.55ex plus .1ex}%
  {\normalfont\large\bfseries}}
\renewcommand{\@maketitle}{%
  \newpage\null\vskip -1.2em%
  {\footnotesize\bfseries
    \href{https://www.ufl.edu/}{University of Florida}, \href{https://math.ufl.edu/}{Department of Mathematics}\par}%
  \vskip 0.18em%
  \tagstructbegin{tag=Title}%
    {\LARGE\bfseries\href{https://gma.math.ufl.edu/exams/algebra-first-year/}{Algebra First-Year Exam}, May 2024, Part 2\par}%
  \tagstructend%
  \vskip 0.35em%
  \tagmcbegin{artifact}%
    \hrule height 1pt%
  \tagmcend%
  \vskip 0.55em%
}
\makeatother
\title{Algebra First-Year Exam, May 2024, Part 2}
\author{}
\date{}
\begin{document}
\maketitle
\thispagestyle{empty}
\begin{examinstructions}
Answer four problems. (If you turn in more, the first four will be graded.) Within reason, you may quote theorems as long as you state them clearly.
\end{examinstructions}
\begin{examproblems}{0}
\item
Prove the following version of the Chinese remainder theorem: Let~\(R\) be a commutative ring with~\(1\) and let \(I, J\) be ideals of~\(R\) such that \(I+J=R\). Then \(IJ=I\cap J\) and there is a ring isomorphism \(R/IJ\to (R/I)\times(R/J)\).
\item
Let~\(R\) be a commutative ring and let~\(M\) be a maximal ideal of~\(R\).
\begin{examsubparts}
\item[\textnormal{(a)}]
Prove that if~\(R\) has an identity then \(R/M\) is a field.
\item[\textnormal{(b)}]
Give an example of a commutative ring~\(R\) and a maximal ideal~\(M\) of~\(R\) such that \(R/M\) is not a field.
\end{examsubparts}
\item
Let~\(R\) be a commutative ring with identity~\(1\) different from~\(0\).
\begin{examsubparts}
\item[\textnormal{(a)}]
Define what is a prime ideal of~\(R\).
\item[\textnormal{(b)}]
Define what is a minimal prime ideal of~\(R\).
\item[\textnormal{(c)}]
Prove that~\(R\) has a minimal prime ideal.
\end{examsubparts}
\item
Prove that the ring \(\mathbb{Z}[\sqrt{-2}]=\{a+b\sqrt{-2}:a,b\in\mathbb{Z}\}\) is a Euclidean domain with respect to the norm defined by \(N(a+b\sqrt{-2})=a^2+2b^2\) for all \(a,b\in\mathbb{Z}\).
\item
Determine the conjugacy classes of elements of order~\(8\) in \(GL(5,\mathbb{Q})\), and, for each conjugacy class, give a representative matrix. (You may use the fact that \(x^4+1\) is irreducible in \(\mathbb{Q}[x]\).)
\item
Let \(F=\mathbb{Q}(\sqrt[4]{2},\sqrt{-1})\). Calculate the degree of the extension \(F/\mathbb{Q}\), and justify each of the steps of your calculation.
\end{examproblems}
\end{document}
